\documentclass[10pt,a4paper]{amsart}
\usepackage[margin=19mm]{geometry}
\usepackage{amsmath,amssymb,mathtools}
\usepackage[scr=rsfs,cal=euler]{mathalfa}
\usepackage{xfrac}
\usepackage[unicode,hidelinks,bookmarks=false]{hyperref}
\newcommand{\ie}{i.e.\ }
\newcommand{\cf}{cf.\ }
\newcommand{\resp}{respectively\ }
\newcommand{\Subsection}{\subsection}
\providecommand{\nobreakdash}{\nobreak\hbox{-}}
\setlength{\emergencystretch}{2em}
\multlinegap0pt
\usepackage{mathtools}
\usepackage{bm}
\usepackage{xcolor}
\usepackage{amssymb}
\usepackage{algorithm}\usepackage{algpseudocode}
\usepackage{float}
\usepackage{multirow}
\usepackage{natbib}
\usepackage{booktabs}
\usepackage{nicefrac}
\newtheorem{lemma}{Lemma}
\newcommand{\dist}{\operatorname{dist}}
\title{Data Completion for Electrical Impedance Tomography by Conditional Diffusion Models}
\author{Ke Chen, Haizhao Yang, Chugang Yi}
\date{Source: arXiv:2602.07813v1 (CC BY 4.0)}
\begin{document}
\maketitle

\noindent\emph{Source and licence.} Data Completion for Electrical Impedance Tomography by Conditional Diffusion Models. Version arXiv:2602.07813v1 (CC BY 4.0), distributed by arXiv under the Creative Commons Attribution 4.0 International licence, http://creativecommons.org/licenses/by/4.0/, as recorded in the arXiv licence metadata for this exact version and shown on the abstract page https://arxiv.org/abs/2602.07813v1. Full licence notice: \texttt{licenses/arxiv-2602.07813-CC-BY-4.0.txt}.\par\smallskip

\noindent\textit{Editorial scope.} Counted unit: the article's lemma lem:rho\_s (source0.tex lines 958--962). The article's complete original proof of it is delivered with it (source0.tex lines 1875--1938, one \texttt{proof} environment of 64 lines, which the article places away from the statement). No mathematics is written, repaired or re-derived: the statement and the proof are the article's own lines, in the article's own notation, and no retained line is substituted or re-flowed.  Retained uncounted as the source-local context the proof needs: lines 1825--1835.  Nothing was omitted from any retained span, so the package carries no omission record.  No retained line contains a citation, so the wrapper carries no bibliography and no bibliography entry.  No retained line contains a figure environment, an image-inclusion command or drawing code, so no image is delivered and the package declares no asset dependency.  Editorial formatting only, in addition: an AMS \texttt{amsart} preamble that restates the article's own declarations and macros used by the retained lines, namely the environments \texttt{lemma} on separate counters, together with the standard packages those lines need.  The retained environments print in this document as Lemma 1 in order of appearance, whereas the article prints the same items with the numbering of its own full document.  The complete native source of the article as distributed in the arXiv e-print for this exact version is bundled unchanged as \texttt{sources/194213/source0.tex}.
\begin{lemma}[Boundary perturbation under vertex perturbations]\label{lem:dist-boundary-perturb}
Let $v,w\in \mathcal{V}^{n_{\mathrm v}}$ be counterclockwise ordered vertex, then
\[
d_H\big(\partial P(v),\partial P(w)\big)\ \le\ \|v-w\|_{\mathrm{poly}}.
\]
Moreover, if $P(v)\subset\Omega$ and $P(w)\subset\Omega$, then
$\dist(P(u),\partial\Omega)=\dist(\partial P(u),\partial\Omega)$ for $u=v,w$\,, and
\[
\dist\big(P(v),\partial\Omega\big)\ \ge\ \dist\big(P(w),\partial\Omega\big)\ -\ \|v-w\|_{\mathrm{poly}}.
\]
\end{lemma}

\begin{lemma}[Local admissibility radius]\label{lem:rho_s}
Fix $P^\ast\in\mathcal A_{1/2}$ and let $v^\ast$ be its counterclockwise ordered vertex list.
Then there exists $\rho_{P^\ast}>0$ such that whenever $\|v-v^\ast\|_{\mathrm{poly}}<\rho_{P^\ast}$,
the polygonal chain $\partial P(v)$ is simple and $P(v)\in\mathcal A_{1/2}$.
\end{lemma}

\begin{proof}[Proof of Lemma ~\ref{lem:rho_s}]
Write $e_i^\ast\coloneqq [v_i^\ast,v_{i+1}^\ast]$ for the edges of $\partial P^\ast$.
Since $P^\ast$ is convex, $\partial P^\ast$ is simple, and every pair of nonadjacent edges is disjoint.
Hence
\[
\eta^\ast\coloneqq \tfrac12\min\big\{\dist(e_i^\ast,e_j^\ast):\ e_i^\ast,e_j^\ast\ \text{nonadjacent}\big\}>0.
\]
Moreover, since $P^\ast\in\mathcal A_{1/2}$, we have the strict margins
\[
m_0\coloneqq \dist(P^\ast,\partial\Omega)-\tfrac12 d_0>0,\qquad
m_1\coloneqq \min_i |v_{i+1}^\ast-v_i^\ast|-\tfrac12 d_1>0,
\qquad
\beta_s\coloneqq \min_i\Big\{\beta_i(v^\ast)-\tfrac12\beta_0,\ \pi-\tfrac12\beta_0-\beta_i(v^\ast)\Big\}>0.
\]

For any $i$,
\[
\big||v_{i+1}-v_i|-|v_{i+1}^\ast-v_i^\ast|\big|
\le |(v_{i+1}-v_{i+1}^\ast)-(v_i-v_i^\ast)|
\le 2\|v-v^\ast\|_{\mathrm{poly}}.
\]
Thus if $\|v-v^\ast\|_{\mathrm{poly}}<m_1/4$, then $|v_{i+1}-v_i|>d_1/2$ for all $i$, i.e.\ \textup{(A2$'$)} holds.


For each edge $e_i(v)=[v_i,v_{i+1}]$, the same interpolation argument as in Lemma~\ref{lem:dist-boundary-perturb}
gives $d_H(e_i(v),e_i^\ast)\le \|v-v^\ast\|_{\mathrm{poly}}$.
Hence for any nonadjacent edge index pair $(i,j)$,
\[
\dist(e_i(v),e_j(v))
\ \ge\ \dist(e_i^\ast,e_j^\ast)-d_H(e_i(v),e_i^\ast)-d_H(e_j(v),e_j^\ast)
\ \ge\ 2\eta^\ast-2\|v-v^\ast\|_{\mathrm{poly}}.
\]
If $\|v-v^\ast\|_{\mathrm{poly}}<\eta^\ast$, then $\dist(e_i(v),e_j(v))>0$ for all nonadjacent pairs,
so nonadjacent edges cannot intersect. Moreover, once we also enforce the angle lower bound $\beta_i(v)>\beta_0/2$ for all $i$ (proved below),
adjacent edges cannot overlap. Hence $\partial P(v)$ is a simple closed polygonal curve.


Now we show the inclusion in $\Omega$ and the boundary margin.
Since $P^\ast\subset\Omega$ and $\dist(P^\ast,\partial\Omega)>d_0/2$, every vertex satisfies
$\dist(v_i^\ast,\partial\Omega)\ge \dist(P^\ast,\partial\Omega)>0$.
Set
\[
\eta_\Omega^\ast\coloneqq \tfrac12\,\dist(P^\ast,\partial\Omega)>0.
\]
If $\|v-v^\ast\|_{\mathrm{poly}}<\eta_\Omega^\ast$, then $|v_i-v_i^\ast|<\eta_\Omega^\ast$ for all $i$,
hence $v_i\in\Omega$ for all $i$. As $\Omega$ is convex, each edge $[v_i,v_{i+1}]\subset\Omega$ and therefore
$P(v)\subset\Omega$. Furthermore, by Lemma~\ref{lem:dist-boundary-perturb},
\[
\dist(P(v),\partial\Omega)\ge \dist(P^\ast,\partial\Omega)-\|v-v^\ast\|_{\mathrm{poly}}.
\]
Thus if $\|v-v^\ast\|_{\mathrm{poly}}<m_0$, then $\dist(P(v),\partial\Omega)>d_0/2$, i.e.\ \textup{(A1$'$)} holds.

We now consider the angle bounds and convexity in (A3$'$). On the set where \textup{(A2$'$)} holds, the interior angles $\beta_i(v)$ are continuous functions of $v$
(e.g.\ they can be expressed via dot/cross products of the adjacent edge vectors).
Hence there exists $\rho_{\mathrm{ang}}>0$ such that if $\|v-v^\ast\|_{\mathrm{poly}}<\rho_{\mathrm{ang}}$ then
$|\beta_i(v)-\beta_i(v^\ast)|<\beta_s/2$ for all $i$, and consequently
$\beta_0/2<\beta_i(v)<\pi-\beta_0/2$ for all $i$, which yields \textup{(A3$'$)}.

Finally, define
\[
\rho_{P^\ast}\coloneqq \min\Big\{\eta^\ast,\ \eta_\Omega^\ast,\ m_0,\ m_1/4,\ \rho_{\mathrm{ang}}\Big\}.
\]
Then $\|v-v^\ast\|_{\mathrm{poly}}<\rho_{P^\ast}$ implies \textup{(A1$'$)--(A3$'$)}.
\end{proof}

\end{document}
